\documentclass[11pt]{article}
\usepackage[a4paper,margin=1in]{geometry}
\usepackage{amsmath,amssymb,mathtools,bm}
\usepackage{enumitem,booktabs,array}
\usepackage{tcolorbox}
\usepackage{hyperref}
\usepackage[protrusion=true,expansion=false]{microtype}
\hypersetup{colorlinks=true,linkcolor=blue,urlcolor=blue,citecolor=blue}
\setlist[itemize]{topsep=2pt,itemsep=2pt}
\setlist[enumerate]{topsep=2pt,itemsep=2pt}
\newtcolorbox{keybox}{colback=blue!3!white,colframe=blue!40!black,boxrule=0.4pt,arc=1mm}
\newtcolorbox{alertbox}{colback=red!3!white,colframe=red!40!black,boxrule=0.4pt,arc=1mm}
\newtcolorbox{answerbox}{colback=green!3!white,colframe=green!40!black,boxrule=0.4pt,arc=1mm}
\newtcolorbox{drillbox}{colback=yellow!5!white,colframe=orange!60!black,boxrule=0.4pt,arc=1mm}
\newcommand{\EV}{\mathbb{E}}
\newcommand{\Var}{\mathrm{Var}}
\newcommand{\Cov}{\mathrm{Cov}}
\newcommand{\blank}[1]{\underline{\hspace{#1}}}

\title{E02-04: Lewellen \& Lowry (2021, JFE) \\
\large ``Does common ownership really increase firm coordination?'' --- Active Recall Workbook}
\author{Empirical Corporate Finance II --- Exam Prep}
\date{\today}
\begin{document}\maketitle

\begin{keybox}
\textbf{Paper in one sentence.} LL directly test the behavioral mechanisms that AST and Backus-Conlon-Sinkinson claim link common ownership to anticompetitive outcomes --- voting, engagement, compensation design, industry-coordination proxies --- and find little evidence that common owners actively coordinate portfolio firms, casting doubt on the supposed mechanism behind the AST price findings.

\textbf{Placement.} Key counterweight to Azar-Schmalz-Tecu (2018). Together they define the common-ownership debate: AST document the correlation with prices; LL finds no behavioral smoking gun for how institutions could cause it.
\end{keybox}

\section*{Round 1: Complete Walk-Through}

\subsection*{1.1 Research question}
AST (2018) propose that common institutional owners cause softer competition by inducing managers to internalize rival firms' profits. What would owners need to \emph{do} to implement this? LL test whether they do it.

\subsection*{1.2 Behavioral channels tested}
\begin{itemize}
\item \textbf{Voting.} Do common owners vote differently on compensation, M\&A, or strategy in ways that dampen competition?
\item \textbf{Engagement / meetings.} Are common owners more likely to meet with or pressure managers toward industry-wide objectives?
\item \textbf{CEO compensation.} Are pay structures at commonly-owned firms more tied to industry profits than to relative performance?
\item \textbf{Industry coordination proxies.} Do common owners reduce product-market rivalry at firms where they hold rivals?
\end{itemize}

\subsection*{1.3 Data and sample}
\begin{itemize}
\item 13F holdings 1980--2017 (ensuring large coverage).
\item ISS voting data, CEO-compensation data, conference-call transcripts, engagement events.
\item Ownership overlap measured consistently with AST.
\end{itemize}

\subsection*{1.4 Headline findings}
\begin{itemize}
\item No robust evidence that common owners vote differently in ways that would soften competition.
\item Engagement activity is dominated by governance (not industry-coordination) themes.
\item CEO pay at commonly-owned firms shows \emph{more} relative-performance evaluation, not less --- opposite to mechanism prediction.
\item Net: behavioral evidence does not support the anticompetitive-coordination story.
\end{itemize}

\subsection*{1.5 Reconciling with AST}
\begin{itemize}
\item Price-level correlation may reflect unobserved heterogeneity.
\item Or may reflect investor preferences but without active coordination.
\item LL suggest the AST causal mechanism is weakly supported, so antitrust remedies should be evaluated with care.
\end{itemize}

\subsection*{1.6 Caveats}
\begin{alertbox}
\begin{itemize}
\item Null findings: absence of evidence is not evidence of absence.
\item Coordination may happen through subtle, unobserved informal channels.
\item Selection: firms with common ownership may be systematically different in ways LL cannot fully absorb.
\end{itemize}
\end{alertbox}

\section*{Round 2: Level 1 Fill-in}
\begin{enumerate}
\item LL tests the \blank{2cm} of common ownership, not the correlation.
\item Channels: voting, engagement, CEO \blank{1.5cm}, industry coordination.
\item CEO pay at commonly-owned firms: \blank{1cm} RPE (not less).
\item Overall finding: \blank{1cm} support for active coordination.
\item Implication: AST correlation may reflect \blank{2cm} heterogeneity.
\end{enumerate}

\section*{Round 3: Level 2 Prompts}
\begin{enumerate}
\item List the behavioral channels LL tests.
\item Summarize the main findings on each channel.
\item How does LL cast doubt on AST's mechanism without refuting the price correlation?
\item What alternative explanations could reconcile AST prices with LL null findings?
\item Discuss antitrust-policy implications.
\end{enumerate}

\section*{Round 4: Minimal Scaffolding}
From a blank page: (i) question, (ii) channels tested, (iii) main findings, (iv) AST reconciliation, (v) policy.

\section*{Round 5: Blank Page}
Essay: common ownership and coordination --- is there a mechanism?

\vspace{6pt}
\begin{drillbox}
\textbf{Speed Drill.} (1) AST claim. (2) LL test approach. (3) Four channels. (4) Main null. (5) Policy implication.
\end{drillbox}

\section*{Quick Reference \& Answer Key}
\begin{answerbox}
\begin{itemize}
\item \textbf{Target.} AST 2018 mechanism.
\item \textbf{Channels tested.} Voting, engagement, pay, coordination.
\item \textbf{Finding.} No behavioral support for active coordination.
\item \textbf{Implication.} AST correlation $\ne$ active coordination.
\end{itemize}
\end{answerbox}

\begin{keybox}
\textbf{30-second exam pitch.} Lewellen \& Lowry (2021) directly interrogate the behavioral mechanism proposed by Azar-Schmalz-Tecu (2018) to explain the anticompetitive effects of common institutional ownership. Using 13F holdings 1980--2017 combined with voting, engagement, compensation, and industry-coordination proxies, they test whether common owners actually do the things the mechanism requires. The behavioral evidence is null or opposite: voting patterns do not favor softer competition, engagement is dominated by governance rather than industry themes, and CEO pay at commonly-owned firms shows \emph{more} relative-performance evaluation (opposite prediction). LL concludes that the AST price-level correlation, if causal, operates through subtler or unobserved channels --- or may reflect omitted variables. The paper anchors the skeptical side of the common-ownership debate and motivates more direct empirical searches for (or against) coordination mechanisms, and it cautions antitrust regulators against acting on AST correlations alone.
\end{keybox}

\end{document}
